\begin{lemma}
Let $F,G\colon\mathbb{C}\to\mathbb{C}$, let $\xi\in\mathbb{C}$, and suppose
that $G$ is real-differentiable at $\xi$ and that $F$ is real-differentiable at
$G(\xi)$.  Then
\[
  \frac{\partial(F\circ G)}{\partial\xi}(\xi)
  =\frac{\partial F}{\partial z}\bigl(G(\xi)\bigr)\frac{\partial G}{\partial\xi}(\xi)
   +\frac{\partial F}{\partial\bar z}\bigl(G(\xi)\bigr)\,
    \overline{\frac{\partial G}{\partial\bar\xi}(\xi)},
\]
\[
  \frac{\partial(F\circ G)}{\partial\bar\xi}(\xi)
  =\frac{\partial F}{\partial z}\bigl(G(\xi)\bigr)\frac{\partial G}{\partial\bar\xi}(\xi)
   +\frac{\partial F}{\partial\bar z}\bigl(G(\xi)\bigr)\,
    \overline{\frac{\partial G}{\partial\xi}(\xi)},
\]
where $\partial/\partial\xi$, $\partial/\partial z$ are as in \noderef{DZ} and
$\partial/\partial\bar\xi$, $\partial/\partial\bar z$ are as in
\noderef{DZbar}.
\end{lemma}
\begin{proof}
For a map $h$ real-differentiable at a point $w$ write $Dh(w)$ for its real
Fr\'echet derivative there, a real-linear map $\mathbb{C}\to\mathbb{C}$.  By
\noderef{DZ} and \noderef{DZbar},
\begin{equation}\tag{0}
  \frac{\partial h}{\partial z}(w)=\tfrac12\bigl(Dh(w)[1]-i\,Dh(w)[i]\bigr),
  \qquad
  \frac{\partial h}{\partial\bar z}(w)=\tfrac12\bigl(Dh(w)[1]+i\,Dh(w)[i]\bigr).
\end{equation}

\emph{Step 1: the real derivative in Wirtinger form.}  Let $h$ be
real-differentiable at $w$.  We claim that
\begin{equation}\tag{1}
  Dh(w)[v]
  =\frac{\partial h}{\partial z}(w)\,v
   +\frac{\partial h}{\partial\bar z}(w)\,\bar v
  \qquad\text{for every }v\in\mathbb{C}.
\end{equation}
Both sides of (1) are $\mathbb{R}$-linear functions of $v$: the left side
because $Dh(w)$ is $\mathbb{R}$-linear, the right side because $v\mapsto v$ and
$v\mapsto\bar v$ are $\mathbb{R}$-linear and the two coefficients are fixed
complex numbers.  Since $1$ and $i$ form a basis of $\mathbb{C}$ over
$\mathbb{R}$, it suffices to verify (1) for $v=1$ and $v=i$.

For $v=1$ we have $\bar v=1$ and, by (0),
\[
  \frac{\partial h}{\partial z}(w)+\frac{\partial h}{\partial\bar z}(w)
  =\tfrac12\bigl(Dh(w)[1]-i\,Dh(w)[i]\bigr)+\tfrac12\bigl(Dh(w)[1]+i\,Dh(w)[i]\bigr)
  =Dh(w)[1].
\]
For $v=i$ we have $\bar v=-i$ and, by (0),
\[
  \frac{\partial h}{\partial z}(w)-\frac{\partial h}{\partial\bar z}(w)
  =\tfrac12\bigl(-2i\,Dh(w)[i]\bigr)=-i\,Dh(w)[i],
\]
so
\[
  \frac{\partial h}{\partial z}(w)\,i+\frac{\partial h}{\partial\bar z}(w)(-i)
  =i\Bigl(\frac{\partial h}{\partial z}(w)-\frac{\partial h}{\partial\bar z}(w)\Bigr)
  =i\bigl(-i\,Dh(w)[i]\bigr)=Dh(w)[i],
\]
using $i(-i)=1$.  This proves (1).

\emph{Step 2: the real chain rule.}  Since $G$ is real-differentiable at $\xi$
and $F$ is real-differentiable at $G(\xi)$, the composite $F\circ G$ is
real-differentiable at $\xi$ and
$D(F\circ G)(\xi)=DF\bigl(G(\xi)\bigr)\circ DG(\xi)$, that is,
\[
  D(F\circ G)(\xi)[v]=DF\bigl(G(\xi)\bigr)\bigl[DG(\xi)[v]\bigr]
  \qquad(v\in\mathbb{C}).
\]
Applying (1) to $h=F$, $w=G(\xi)$ and the vector $DG(\xi)[v]$ gives
\begin{equation}\tag{2}
  D(F\circ G)(\xi)[v]
  =\frac{\partial F}{\partial z}\bigl(G(\xi)\bigr)DG(\xi)[v]
   +\frac{\partial F}{\partial\bar z}\bigl(G(\xi)\bigr)\overline{DG(\xi)[v]} .
\end{equation}

\emph{Step 3: taking the two Wirtinger combinations.}  Abbreviate
$A=\frac{\partial F}{\partial z}(G(\xi))$ and
$B=\frac{\partial F}{\partial\bar z}(G(\xi))$.  By (0) applied to $F\circ G$ at
$\xi$, and then (2) with $v=1$ and $v=i$,
\[
  \frac{\partial(F\circ G)}{\partial\xi}(\xi)
  =\tfrac12\Bigl(D(F\circ G)(\xi)[1]-i\,D(F\circ G)(\xi)[i]\Bigr)
  =A\cdot\tfrac12\bigl(DG(\xi)[1]-i\,DG(\xi)[i]\bigr)
   +B\cdot\tfrac12\Bigl(\overline{DG(\xi)[1]}-i\,\overline{DG(\xi)[i]}\Bigr),
\]
after regrouping the four terms (legitimate since $A$ and $B$ are fixed complex
numbers).  By (0) the first bracket is
$\frac{\partial G}{\partial\xi}(\xi)$.  For the second, $\overline{i\,w}=-i\bar w$
for every $w\in\mathbb{C}$, so $-i\,\overline{DG(\xi)[i]}=\overline{i\,DG(\xi)[i]}$
and, conjugation being additive and $\tfrac12$ real,
\[
  \tfrac12\Bigl(\overline{DG(\xi)[1]}-i\,\overline{DG(\xi)[i]}\Bigr)
  =\overline{\tfrac12\bigl(DG(\xi)[1]+i\,DG(\xi)[i]\bigr)}
  =\overline{\frac{\partial G}{\partial\bar\xi}(\xi)} .
\]
This gives the first identity of the lemma.

For the second, (0) gives
\[
  \frac{\partial(F\circ G)}{\partial\bar\xi}(\xi)
  =\tfrac12\Bigl(D(F\circ G)(\xi)[1]+i\,D(F\circ G)(\xi)[i]\Bigr)
  =A\cdot\tfrac12\bigl(DG(\xi)[1]+i\,DG(\xi)[i]\bigr)
   +B\cdot\tfrac12\Bigl(\overline{DG(\xi)[1]}+i\,\overline{DG(\xi)[i]}\Bigr).
\]
The first bracket is $\frac{\partial G}{\partial\bar\xi}(\xi)$ by (0).  For the
second, $\overline{-i\,w}=i\bar w$, so
$i\,\overline{DG(\xi)[i]}=\overline{-i\,DG(\xi)[i]}$ and
\[
  \tfrac12\Bigl(\overline{DG(\xi)[1]}+i\,\overline{DG(\xi)[i]}\Bigr)
  =\overline{\tfrac12\bigl(DG(\xi)[1]-i\,DG(\xi)[i]\bigr)}
  =\overline{\frac{\partial G}{\partial\xi}(\xi)} .
\]
This gives the second identity.
\end{proof}
